% ============================================================================
\clearpage
\section*{v4.10.0-r35 Source-Only: The Optical Mode of the Preparation as the Mediator of the OP5 Case}
\addcontentsline{toc}{section}{v4.10.0-r35 Source-Only: The Optical Mode as OP5 Mediator}
\label{sec:v4100-r35-material-mediator}
% ============================================================================

\paragraph{Purpose.}
After v4.10.0-r34 no vacuum field of the Companion carries the directed contrast of the OP5 case.  A backward check of the same day shows
why this is forced: a propagation coefficient \(r_m\neq1\) and the pairing of a matter momentum with its counter-propagating tachyon image at a
common rapidity both single out a frame.  Under a boost by \(\xi\) the matter rapidity moves to \(\eta+\xi\), the image to \(\eta-\xi\), and the
squared mass of the pair connector changes by \(4m^2\cosh(2\eta+\xi)\sinh\xi\).  This layer follows the author's choice of the material route:
the frame and the axis come from the preparation, and the mediator is the optical mode of the two-component preparation of r32.2.  It is an
exploration with a method fixed before the computation and is not blind.  Earlier layers are not rewritten.

\noindent\textbf{Source-only status.}
This v4.10.0-r35 layer is source-only.  It does not build a PDF, does not assign a new DOI,
and does not authorize a public release.  The v5.3.0 Version DOI
\href{https://doi.org/10.5281/zenodo.24434403}{10.5281/zenodo.24434403} and the Concept DOI
\href{https://doi.org/10.5281/zenodo.19210728}{10.5281/zenodo.19210728} are unchanged by it.

\subsection*{v4.10.0-r35 results}
\label{subsec:v4100-r35-results}

Throughout, \(x=\omega^2\), \(q=|\mathbf k|^2\), \(a_\perp=|\mathbf k_\perp|^2\), \(\mu=\rho_1\rho_2/(\rho_1+\rho_2)\), and the comparison field is the Higgs radial mode of
W2 (r33.1), coupled to the two components by \(\eta\,(Y_1\nabla\cdot\mathbf u+Y_2\nabla\cdot\mathbf w)\) with \(Y_a=\sum y\,n_0\) of component \(a\) (r32.1).

\paragraph{Theoremlet v4.10.0-r35.1 (isotropic preparation).}
Integrating out the two continua gives the operator \(2(m_H^2+q-x)+q\,Y^{\mathsf T}M^{-1}Y\), where \(M\) is the plane-wave matrix of the two components.  It has an
acoustic pole with residue \(q(Y_1+Y_2)^2/(\rho_1+\rho_2)\) and an optical pole at \(\omega_0^2=k(\rho_1+\rho_2)/(\rho_1\rho_2)\) with residue
\[
 q\,\mu\Bigl(\frac{Y_1}{\rho_1}-\frac{Y_2}{\rho_2}\Bigr)^2+O(q^2).
\]
Both depend on \(q\) only; an isotropic preparation gives no directed term.  \emph{Proof.} Residues of the rational function at the two roots of
\(\det M=0\), expanded in \(q\).  \(\square\)

\paragraph{Theoremlet v4.10.0-r35.2 (the leading static contrast does not use \(r_m\)).}
On the static plateau
\[
 \frac{g_m^2N_m}{d_m}\Big|_{x=0}=\frac{2g_m^2a_\perp^2}{3m^2}-\frac{2g_m^2r_m\,a_\perp^2\,q}{3m^4}+O(q^2),
\]
so the leading term, from which r32.4 obtains the contrast, is independent of \(r_m\).  The subluminal coefficient \(r_m=\tfrac12\) of r32.0 is not
needed for the leading directed contrast.  \emph{Proof.} Series of \(1/(m^2+r_mq)\) in \(q\).  \(\square\)

\paragraph{Theoremlet v4.10.0-r35.3 (uniaxial preparation).}
Let the pair spring be rigid along the axis and equal to \(k_\perp\) across it.  Then the axial component carries only the acoustic pole, and the
transverse optical pole at \(\omega_0^2=k_\perp(\rho_1+\rho_2)/(\rho_1\rho_2)\) has the leading residue
\[
 a_\perp\,\mu\Bigl(\frac{Y_1}{\rho_1}-\frac{Y_2}{\rho_2}\Bigr)^2 ,
\]
which depends on \(a_\perp\) alone.  The statically induced source of the optical part is \(-a_\perp\mu(Y_1/\rho_1-Y_2/\rho_2)^2\eta/\omega_0^2\), the form
\(\Delta_\perp\eta/\omega_0^2\) of the comparison field \(f=\Delta_\perp H/m^2\) of r32.4 with \(m^2\leftrightarrow\omega_0^2\).
\emph{Proof.} With component-wise elastic operators the Cartesian components decouple; the limit of infinite axial spring projects the axial
component onto equal displacements, and the transverse component gives the residue as in r35.1 with \(q\) replaced by \(a_\perp\).  \(\square\)

\paragraph{Remark v4.10.0-r35.4 (reading).}
\begin{center}
\small
\begin{tabular}{@{}>{\raggedright\arraybackslash}p{0.30\textwidth}>{\raggedright\arraybackslash}p{0.62\textwidth}@{}}
\toprule
model input (r32.0) & material counterpart\\
\midrule
mass \(m^2\) & optical gap \(\omega_0^2\)\\
axis, \(P=\operatorname{diag}(1,1,0)\) & uniaxial pair spring\\
coupling \(g_m^2\) & \(\mu(Y_1/\rho_1-Y_2/\rho_2)^2=\mu(f_1-f_2)^2/v^2\) with \(f_a\) the Higgs fraction of the rest mass\\
\(r_m=\tfrac12\) & not needed at leading order (r35.2)\\
numerator \(\tfrac23a_\perp^2\) & \(a_\perp\): one power less\\
\bottomrule
\end{tabular}
\end{center}
Against the criteria fixed beforehand the material route carries the directed contrast \emph{partially}: mass, axis and coupling are properties
of the sample and can be varied independently, the coupling exists only if the two components carry different Higgs fractions of their mass,
and the induced source has the form of \(f\).  The operator of the comparison model contains the back-reaction \(g_m^2\Delta_\perp f\), one power of
\(a_\perp\) more than a density coupling supplies; a coupling to second derivatives of the displacement would be needed and is expected to be
suppressed by the square of lattice spacing times wave number (estimate, not computed).  Which power the OP5 readout measures, and the size of
\(f_1-f_2\), stay open.

\subsection*{v4.10.0-r35 frames}
\label{subsec:v4100-r35-frames}

All quantities are computed in the matter frame M; the boost check of the pairing rule compares the two rapidities under the same boost and
no quantity of another frame is used.

\subsection*{v4.10.0-r35 register entries}
\label{subsec:v4100-r35-register}

\begin{center}
\small
\begin{tabular}{@{}>{\raggedright\arraybackslash}p{0.34\textwidth}>{\raggedright\arraybackslash}p{0.58\textwidth}@{}}
\toprule
statement & status after v4.10.0-r35\\
\midrule
directed mediator needs a distinguished frame & shown for \(r_m\neq1\) and for the pair connector (purpose)\\
optical mode of a uniaxial preparation & directed, massive, coupled iff \(f_1\neq f_2\); numerator one power short\\
\(r_m=\tfrac12\) in the leading contrast & not needed (r35.2)\\
\bottomrule
\end{tabular}
\end{center}

\subsection*{v4.10.0-r35 notes}
\label{subsec:v4100-r35-notes}

\texttt{notes/\allowbreak research/\allowbreak op5\_\allowbreak material\_\allowbreak mediator\_\allowbreak 2026\_\allowbreak 11\_\allowbreak 22/} (method note commit a86479e0,
script \texttt{material\_\allowbreak mediator.wl}, output in \texttt{run/}).

\subsection*{v4.10.0-r35 ledger and guard update}
\label{subsec:v4100-r35-ledger-guard}

\[
R_{\rm chk}^{v4100\text{-}r35}=(1,1,0)
\]
for (status recorded; computations checked in the cited script; the second-gradient coupling and the readout power not derived).

The layer spends no guard:
\[
R_{\rm glob}^{(v4100\text{-}r35,+)}=(0,1,1,1,1),
\qquad
\Gamma_{\rm guard}^{v4100\text{-}r35}=(1,1,1,1,1),
\qquad
\beta_4=1/256\ \text{unspent}.
\]
The layer does not imply exact finite \(C5\) zero closure, exact all-readout terminality, \(C6/T2\), a global photon-clock/RH theorem,
or global all-readout closure.  OP1--OP5, C5 and C6/T2 remain open.
